\documentclass{article}
\usepackage[T1]{fontenc}
%\usepackage{gensymb}
%\usepackage{microtype}
\usepackage[margin=3cm]{geometry}
\usepackage{lmodern}
\usepackage{amsmath}
\usepackage{amssymb}
\usepackage{titlesec}
\usepackage{enumitem}
% {left-indent}{space before}{space after}
\titlespacing*{\subsection}{0pt}{2.5em plus 1ex minus .2ex}{2ex plus .2ex}
\titlespacing*{\section}{0pt}{6ex plus 1ex minus .2ex}{2.3ex plus .2ex}
\linespread{1.2}
\begin{document}

\section{Bijection}
\subsection{Bijection Problem}
Consider $\ln : (0,\infty) \to \mathbb{R}$. Is it injective? Is it surjective? Is it bijective?
\begin{enumerate}
  \item Let $f(x) = e^{x}$ and $g(x) = \ln(x)$
    \begin{align*}
      (f \circ g) &= z \\
      z       &= e^{\ln(x)} \\
      \ln(z)  &= \ln(x) = t && \text{(adding var t)} \\
      e^{t}      &= z = x \\
      z       &= x \\
      (g \circ f) &= y \\
      y       &= \ln(e^{x}) \\
      e^{y}      &= e^{x} \\
      y       &= x
    \end{align*}
\end{enumerate}

\newpage
\subsection{Second Attempt}
Proove $e^{x}$ is injective where $(0, \infty) \to \mathbb{R}$: \\
$w_{2}>w_{1}$ and $w_{2}-w_{1} > 0$
\begin{align*}
  e^{w_{2}} &= e^{w_{1}} \\
  e^{w_{2}-w_{1}} &= 1 \\
  w_{2} - w_{1} &= 0 \\
  \implies w_{2} &= w_{1} \qquad \text{contradiction}
\end{align*} 
Proove: $\ln(a) = \ln(b)$
\begin{align*}
  \ln(a) &= t \\
  a      &= e^{t} \\
  \ln(b) &= t \\
  b      &= e^{t} \\
  a      &= b
\end{align*}
Proove $\ln(x) = y$ is surjective:
\begin{align*}
  \ln(x) &= y \\
  \implies e^{y} &= x \\
\end{align*}

\newpage
\subsection{Questions}
\begin{enumerate}
  % Question 1:
  \item Proof: Let $a$, $b \in \mathbb{R}_{>0}$ such that:
    \begin{equation}
      \log_{2}(a) = \log_{2}(b)
    \end{equation}
    Raise 2 to the power of each side:
    \begin{equation}
      2^{\log_{2}(a)} = 2^{\log_{2}(b)}
    \end{equation}
    Since $2^{x}$ and $\log_{2}(x)$ are inverse functions:
    \begin{equation}
      a = b
    \end{equation}
    since $a = b$, the function $\log_{2}$ is injective. This is provable without previous proofs.
  
  \setcounter{equation}{0}
  % Question 2:
  \item Is $e^{x} : \mathbb{R} \to \mathbb{R}$ surjective: \\
    Let $y = e^{x}$. For any value of $y$ where $y \le 0$, so 0 or negative. No value of x can make 
    $e^{x}$ 0 or negative, so the codomain has to be bound to $\mathbb{R}_{>0}$ for $e^{x}$ to be surjective.
    Solution: 
    \begin{equation}
      e^{x} : \mathbb{R} \to \mathbb{R}_{>0} 
    \end{equation}
    
    \setcounter{equation}{0}
  % Q3:
  \item $x^{2} : \mathbb{R} \to \mathbb{R}$ is not injective, as for any $x$, $-x$ produces the same output.
    Restricting the domain $\mathbb{R}_{\ge 0}$ fixes this. The function also can't output negatives even without a restricted domain so:
    \begin{equation}
      x^{2} : \mathbb{R}_{\ge 0} \to \mathbb{R}_{\ge 0} 
    \end{equation}
    0 is included because:
    \begin{equation}
      0^{2} = 0 
    \end{equation}
    So it maps correctly, and one to one. So for this restricted domain and codomain, $x^{2}$ is injective. 
    The inverse of a square is a square root where:
    \begin{equation}
      \sqrt x^{2} = x
    \end{equation}
    And the $\sqrt x$ is valid and injective for any $\mathbb{R}_{\ge 0} \to \mathbb{R}_{\ge 0}$,
    but no real values exist for negatives.
  
  \setcounter{equation}{0}
  \newpage
  %Q4:
  \item $x^{3} : \mathbb{R} \to \mathbb{R}$ is bijective because there is an output for each input, including negatives,
    as odd numbered exponentials always have a leftover negative:
    \begin{equation}
      x^{3} = x \cdot x \cdot x
    \end{equation}
    sub -1 for x:
    \begin{equation}
      -1 \cdot -1 \cdot -1 = -1
    \end{equation}
    or -5:
    \begin{equation}
      -5 \cdot -5 \cdot -5 = -125
    \end{equation}
    This prooves the function is surjective, as all positive and negative values map and continue to do so as $x \to \pm \infty$. \\
    To prove injectivity, let $a^{3} = b^{3}$ where $b \neq 0$
    \begin{equation}
      \implies \left(\frac{a}{b} \right)^{3} = 1
    \end{equation}
    \begin{equation}
      \implies \frac{a}{b} = 1
    \end{equation}
    \begin{equation}
      a = b
    \end{equation}
    Because $0^{x}$:
    \begin{equation}
    a^{3} \implies a = 0 = b  
    \end{equation}
    Therefore $x^{3} : \mathbb{R} \to \mathbb{R}$ is injective. And as it is also surjective, it is therefore bijective. 
      
\end{enumerate}


\end{document}
